% related_work_patch.tex -- proposed changes to Section 2 (Related work) of main.tex.
% Round-2 literature review, 12 November 2026. Not \input anywhere; copy by hand.
% New keys are in review/references_r2.bib; existing keys are from references.bib.
% Each block says whether it REPLACES or is ADDED to an existing paragraph.

%----------------------------------------------------------------------
% (1) REPLACE the paragraph "Aggregate flexibility." with:
%----------------------------------------------------------------------
\paragraph{Aggregate flexibility.}
For a fixed EV population, aggregate flexibility is a Minkowski sum of
individual sets \citep{mukhiAbate2023exact,altaha2024flexibility}.
\citet[Thms.~3 and~7]{mukhiLohoAbate2025exact} characterize that sum and
disaggregate its schedules for generalized polymatroids;
\citet{pandaTindemans2024} give a representation whose size depends only on
the number of time steps, and majorization-type conditions characterize
fleets of rate-limited loads \citep{nayyar2016duration}.
\citet[Eqs.~(15)--(17)]{muller2019aggregation} also compute the least cost of
disaggregation for zonotopic resources with structured costs.
\citet{arvis2026aggregated} embed aggregate EV flexibility in unit commitment,
and \citet{altahaBitar2024lossy} give a state-of-charge reformulation for a
single lossy storage device. All of these take the population of vehicle
tasks as given. In our setting a matching determines which duties enter the
sum. For lossless charging only, a fixed matching's aggregate is an interval
and is covered by the results above; the content of
\cref{thm:projection} is that, at a fixed number of pairs, one matching's
set \emph{contains} every other matching's, and that this survives
conversion losses, exports and a shared gross connector budget, where the
net-interval aggregate is no longer exact (\cref{rem:netinterval}).
Computing the representative is a minimum-cost $k$-cardinality assignment
\citep{dellamicoMartello1997kcard}; its value is convex in $k$ as a
minimum-cost flow value \citep{murota2003dca}, and for the complete graph
the opposite-order pairing follows from the Monge property of
$(a_i+b_j-B)_+$ \citep{burkard1996monge}.

%----------------------------------------------------------------------
% (2) REPLACE the last two sentences of "Nonconvex pricing."
%     ("Our gap bound uses the distance ... \citep{starr1969,hreinsson2021}.")
%     with:
%----------------------------------------------------------------------
Shapley--Folkman-type bounds control the gap between a separable
nonconvex problem and its convexification by the number of coupling
constraints times the largest individual nonconvexity
\citep{starr1969,aubinEkeland1976,bertsekas1983scheduling,udellBoyd2016gap};
\citet{hreinsson2021} apply this to dispatchable demand, and
\citet{chao2019incentives} shows that side payments become negligible
relative to surplus as markets grow. Applied to a fleet with buses or
services as blocks, these bounds scale with the number of services, because
the matching constraints couple the blocks. Our hull convexifies whole-fleet
plans, so the gap is at most the nonconvexity of a single fleet value
function. \Cref{prop:onefee} shows that this nonconvexity is below one
bus-day fee. The argument is the convexity of minimum-cost matching in its
cardinality; equivalently, an optimum of the matching polytope with one side
constraint lies between two adjacent matchings, whose sizes differ by at
most one \citep{raviGoemans1996cmst,berger2011budgeted}. We use these
classical facts rather than claim a new duality-gap technique.

%----------------------------------------------------------------------
% (3) ADD after the paragraph "Price adjustment." (or merge into it):
%----------------------------------------------------------------------
\paragraph{Coupled power and transport prices.}
Price-response loops between electric mobility and the grid are known to
cycle. In \citet[Sec.~V-A and Table~II]{alizadeh2017coupled}, a power
operator that prices the previous period's charging pattern drives traffic
and charging into a two-period alternation, although the coupled problem is
convex and has a socially optimal price. \citet{wei2018equilibrium} compute
equilibria of a traffic-assignment and distribution-OPF loop. They establish
existence under continuity and only local, initialization-dependent
convergence. \citet{roozbehani2012volatility} give
elasticity conditions under which real-time pricing feedback is unstable.
Distribution prices can decentralize divisible EV charging at a
clearing price \citep{liWuOren2014dlmp}. Proximal or damped updates restore
convergence for continuous charging
\citep{ganTopcuLow2013,ma2016decentralized}, and set-valued LMP updates
converge to equilibria of divisible flexible loads \citep{gong2019lmp}.
Closer to buses, \citet{yetkin2024cooptimizing} co-optimize transit V2G and
DC dispatch with nodal prices, and \citet{najafi2025joint} couple binary
trip chains to a conic distribution model after fixing the cover.
A price cycle therefore says little on its own. \Cref{thm:iteration} adds
that, with an indivisible whole-fleet response, undamped updates cycle
whenever the hull optimum lies inside a face, with or without a support gap. Only the face's
geometry, not the cycle, reveals whether a supporting price exists.

%----------------------------------------------------------------------
% (4) ADD at the end of "Electric-bus scheduling with charging." :
%----------------------------------------------------------------------
Whether duties can be fixed before charging is optimized has so far been
answered computationally. \citet{sassiOulamara2017} assign tours first and
schedule charging as a minimum-cost flow. Reviews treat vehicle and charge
scheduling as separate but linked problems \citep{perumal2022review}.
On 14 transit networks, \citet{nath2024joint} report that integrating
charge scheduling into trip assignment lowers total cost by 4.1\% on average
relative to a sequential plan, including charger siting. In logistics,
flexible departures save up to about 5\% under day-ahead prices
\citep{kleinSchiffer2023flexible}. Much larger savings arise from optimizing
charging on a fixed schedule: 16.5\% against greedy rules under day-ahead
prices \citep{abdelwahed2020opportunity}, and more where demand charges
dominate \citep{heLiuSong2020}. They also arise from choosing fleet size
jointly with time-of-use-aware charging \citep{wu2022grid,zhang2024tou}.
We are not aware of a result stating conditions under which a
price-independent duty choice loses nothing. \Cref{thm:projection} gives
such conditions for a narrow class, and the boundary test in
\cref{sec:boundaries} (at most about 2\%) is in line with the measured
sequential-versus-integrated gaps.

%----------------------------------------------------------------------
% (5) REPLACE the clause on Zhou et al. in "Buses, tariffs and markets."
%     ("...and \citet{zhou2025v2g} formulate ... published abstract).")
%     with (full text now read, Kyoto repository accepted manuscript):
%----------------------------------------------------------------------
and \citet{zhou2025v2g} jointly choose fleet size, trip chains and
charging and discharging for 418 trips at two depots under a three-level
time-of-use tariff, solved by adaptive large-neighbourhood search; in their
case cross-line operation lowers daily cost by about 17\% relative to
single-line operation (their Table~6).
